\documentclass[12pt,a4paper,twoside,openany]{book}

% Packages
\usepackage[utf8]{inputenc}
\usepackage[T1]{fontenc}
\usepackage[english]{babel}
\usepackage{amsmath,amssymb,amsthm}
\usepackage{graphicx}
\usepackage{subcaption}
\usepackage{float}
\usepackage{hyperref}
\usepackage{cite}
\usepackage{algorithm}
\usepackage{algorithmic}
\usepackage{listings}
\usepackage{xcolor}
\usepackage{geometry}
\usepackage{fancyhdr}
\usepackage{setspace}
\usepackage{tikz}
\usetikzlibrary{shapes,arrows,positioning,fit,calc}

% Page geometry
\geometry{
    left=3cm,
    right=2.5cm,
    top=2.5cm,
    bottom=2.5cm
}

% Header and footer
\pagestyle{fancy}
\fancyhf{}
\fancyhead[LE,RO]{\thepage}
\fancyhead[RE]{\leftmark}
\fancyhead[LO]{\rightmark}
\setlength{\headheight}{14.5pt}

% Code listings settings
\lstset{
    basicstyle=\ttfamily\footnotesize,
    keywordstyle=\color{blue},
    commentstyle=\color{green!60!black},
    stringstyle=\color{red},
    numbers=left,
    numberstyle=\tiny\color{gray},
    stepnumber=1,
    numbersep=8pt,
    showstringspaces=false,
    breaklines=true,
    frame=single,
    backgroundcolor=\color{gray!10}
}

% WebAssembly language definition
\lstdefinelanguage{wasm}{
    keywords={func, loop, block, if, else, end, local, global, param, result, return, call, call_indirect, drop, select, get_local, set_local, tee_local, get_global, set_global, i32, i64, f32, f64, load, store, const, eqz, eq, ne, lt, gt, le, ge, add, sub, mul, div, and, or, xor, shl, shr, rotl, rotr, clz, ctz, popcnt, br, br_if, br_table, unreachable, nop, memory, table, elem, data, offset, align, import, export, start, type, module},
    morekeywords={[2]suspend, resume, cont.new, cont.bind, local.get, local.set, local.tee, global.get, global.set, ref.func, ref.null, ref.is_null},
    keywordstyle=\color{blue}\bfseries,
    keywordstyle={[2]\color{purple}\bfseries},
    comment=[l]{;;},
    morecomment=[s]{(;}{;)},
    commentstyle=\color{green!60!black}\itshape,
    morestring=[b]",
    stringstyle=\color{red},
    sensitive=true,
    alsoletter={.\$},
    literate=*{(}{{{\color{gray}(}}}1 {)}{{{\color{gray})}}}1 {\$}{{{\color{orange}\$}}}1
}

% Kotlin language definition
\lstdefinelanguage{Kotlin}{
    keywords={abstract, actual, as, as?, break, by, class, companion, continue, data, do, dynamic, else, enum, expect, false, final, for, fun, get, if, import, in, infix, inline, interface, internal, is, lateinit, null, object, open, operator, out, override, package, private, protected, public, return, sealed, set, super, suspend, this, throw, true, try, typealias, typeof, val, var, vararg, when, where, while},
    morekeywords={[2]Any, Boolean, Byte, Char, Continuation, CoroutineContext, Double, EmptyCoroutineContext, Float, Int, Long, Nothing, Result, Short, String, Unit, COROUTINE_SUSPENDED},
    keywordstyle=\color{blue}\bfseries,
    keywordstyle={[2]\color{purple}\bfseries},
    comment=[l]{//},
    morecomment=[s]{/*}{*/},
    commentstyle=\color{green!60!black}\itshape,
    morestring=[b]",
    stringstyle=\color{red},
    sensitive=true,
    alsoletter={?},
}

% Theorem environments
\newtheorem{theorem}{Theorem}[chapter]
\newtheorem{lemma}[theorem]{Lemma}
\newtheorem{proposition}[theorem]{Proposition}
\newtheorem{corollary}[theorem]{Corollary}
\theoremstyle{definition}
\newtheorem{definition}[theorem]{Definition}
\newtheorem{example}[theorem]{Example}
\theoremstyle{remark}
\newtheorem{remark}[theorem]{Remark}

% Hyperref settings
\hypersetup{
    colorlinks=true,
    linkcolor=blue,
    citecolor=blue,
    urlcolor=blue,
    pdftitle={Computer Science Thesis},
    pdfauthor={Your Name}
}

% Line spacing
\onehalfspacing

% Document information
\title{Your Thesis Title}
\author{Your Name}
\date{\today}

\begin{document}

% Front matter
\frontmatter
\begin{titlepage}
    \thispagestyle{empty}
    \centering
    \vspace*{2cm}
    
    {\LARGE\bfseries Constructor University\par}
    \vspace{0.5cm}
    {\large Department of Computer Science\par}
    
    \vspace{3cm}
    
    {\huge\bfseries Compiling Kotlin Coroutines to WebAssembly with Stack-switching\par}
    
    \vspace{2cm}
    
    {\Large A thesis submitted in partial fulfillment\par}
    {\Large of the requirements for the degree of\par}
    \vspace{0.5cm}
    {\Large\bfseries Master of Science in Advanced Software Technology\par}
    
    \vspace{2cm}
    
    {\large\bfseries Veronika Sirotkina\par}
    
    \vfill
    
    {\large Supervisor: Prof. Dr. Anton Podkopaev\par}
    {\large Industry supervisor: Zalim Bashorov\par}

    \vspace{1cm}

    {\large \today\par}
    
\end{titlepage}

\chapter*{Abstract}
\addcontentsline{toc}{chapter}{Abstract}

Kotlin/Wasm is a compilation target for the Kotlin programming language that compiles Kotlin code to
WebAssembly. It is used, for example, by the Compose Multiplatform library for its Web target.
Kotlin Coroutines are central to Kotlin/Wasm applications, yet the current implementation is based on
state-machine transformation: each suspend function is compiled into a state machine that saves and
restores local variables at each suspension point. This approach adds code that increases binary size
and introduces runtime overhead.

The WebAssembly Stack-switching proposal provides native primitives for suspending and resuming
execution contexts. This thesis presents an alternative implementation of Kotlin Coroutines for
Kotlin/Wasm based on these primitives, replacing the state-machine transformation with direct
stack-switching instructions.

A key technical challenge addressed in this work is the mismatch between Kotlin's coroutine model,
where a continuation object must be available before a function suspends, and the Wasm Stack-switching
model, where the continuation only becomes available after suspension. The thesis describes two
implementations that resolve this mismatch.

The evaluation shows that the Stack-switching implementation reduces binary size by approximately
10.8\% in production builds and 2.3\% in development builds for a small Compose Multiplatform
project. However, performance benchmarks on V8 show a noticeable regression: the most realistic
benchmark cases are around 2 times slower than the baseline, and cases with intensive suspension
chaining are 4--4.5 times slower without an additional V8 flag, or around 1.5 times slower with it.
The performance regression is likely related to the current state of the V8 Stack-switching
implementation, which is still under active development.


\chapter*{Acknowledgments}
\addcontentsline{toc}{chapter}{Acknowledgments}

I would like to thank Zalim Bashorov for guiding me through the initial development process and helping me navigate the
compiler code. I would also like to thank Aleksandr Shefer, who took over the implementation and is now bringing it to
production.

\tableofcontents
\listoffigures
\listoftables
\lstlistoflistings

% Main matter
\mainmatter
\chapter{Introduction}
\label{ch:introduction}

\section{Motivation}
\label{sec:motivation}

WebAssembly is a binary instruction format for a stack-based virtual machine. Its most mainstream use case is in browsers. Many languages can be compiled to it, including Kotlin.

There is a proposal to implement stack-switching in WebAssembly. The specification for stack-switching is finalized and its implementation is on the way.

The current implementation of Kotlin Coroutines for Wasm (with state-machines and continuation passing style for suspend function) is based on the implementation of suspend functions for Kotlin/JVM, which doesn't have primitives for continuation objects and suspension of functions.

The ultimate goal is to verify whether compiling Kotlin Coroutines using stack-switching improves binary size, complexity and performance for available Wasm runtimes with stack-switching support.

There is also an indirect benefit of preparing for eventual release of stack-switching in the mainstream Wasm runtimes.

\section{Thesis Structure}
\label{sec:structure}

The thesis is organized as follows: Chapter~\ref{ch:background} provides background information on WebAssembly stack-switching and Kotlin coroutines primitives. Chapter~\ref{ch:implementation} describes the implementation of Kotlin coroutines using the WebAssembly stack-switching proposal. Chapter~\ref{ch:evaluation} evaluates the implementation in terms of performance and binary size. Finally, Chapter~\ref{ch:conclusion} concludes the thesis and discusses future work.

\chapter{Background and Related Work}
\label{ch:background}

\section{WebAssembly and Stack-switching}
\label{sec:webassembly}

Stack-switching proposal introduces a type for continuation and the corresponding instructions to work with it. Continuation is an object that represents a function at some point in its execution. That point is either the beginning of the function or a marked suspension point. Resuming continuation resumes the function from the recorded point. Continuation can also be intuitively understood as "the rest of a computation."

A more detailed description and an example of a Wasm program with stack-switching can be found in Appendix~\ref{ch:appendix}.

\section{Kotlin Coroutines}
\label{sec:kotlin-coroutines}

Kotlin language offers low-level coroutine primitives and suspending functions. Together these two things are used to implement higher-level concurrency abstractions (most notably, the library kotlinx.coroutines)

There are three notable things in regards to the compilation of Kotlin coroutines:

\begin{enumerate}
    \item Suspend functions have an additional parameter appended to their parameter list. This is a continuation parameter to handle callbacks.
    \item Suspend functions are compiled into a state machine class. Local variables are "spilled" (saved) into this class to preserve the execution state on the heap, allowing the function to resume execution at the correct point after suspension.
    \item Some of the low-level coroutine primitives are provided as intrinsic implementations. This means that they are written directly in the target language instead of being compiled from Kotlin.
\end{enumerate}

\section{Related Work}
\label{sec:related-work}

\subsection{Approach 1}
\label{subsec:approach1}

\subsection{Approach 2}
\label{subsec:approach2}

\subsection{Comparison and Gap Analysis}
\label{subsec:gap-analysis}

\section{Technologies and Tools}
\label{sec:technologies}

\section{Summary}
\label{sec:background-summary}

\chapter{Implementation}
\label{ch:implementation}

\section{Introduction}
\label{sec:implementation-intro}

\section{Development Environment}
\label{sec:dev-environment}

\subsection{Programming Languages and Frameworks}
\label{subsec:languages}

\subsection{Development Tools}
\label{subsec:tools}

\section{System Architecture}
\label{sec:system-architecture}

\section{Core Components}
\label{sec:core-components}

\subsection{Component 1 Implementation}
\label{subsec:impl-component1}

\subsection{Component 2 Implementation}
\label{subsec:impl-component2}

\section{Data Structures and Algorithms}
\label{sec:data-structures}

\section{User Interface}
\label{sec:user-interface}

\section{Testing and Validation}
\label{sec:testing}

\section{Challenges and Solutions}
\label{sec:challenges}

\section{Summary}
\label{sec:implementation-summary}

\chapter{Evaluation}
\label{ch:evaluation}

\section{Correctness}
\label{sec:testing}

Kotlin codebase already contains extensive tests for coroutines. We can separate them into
two levels.

\subsection*{Intrinsic semantics tests}

These tests verify the behavior of the lowest-level coroutine
primitives like \texttt{startCoroutineUninterceptedOrReturn} and \texttt{suspendCoroutineUninterceptedOrReturn}
They also cover low-level coroutine APIs built directly on top of them like \texttt{startCoroutine}
and \texttt{suspendCoroutine}

The tests are located in the Kotlin repository at
\texttt{compiler/testData/codegen/box/coroutines/intrinsicSemantics}

Below is an example of what kind of tests are in the intrinsicSemantics folder

TODO(insert listing from suspendCoroutineUninterceptedOrReturn.kt test)

\subsection*{Box tests for coroutines}

This is an extensive collection of tests that verifies that coroutines work with other Kotlin features.
TODO(check what tests there are and describe 2-3 things that are tested that are easy to understand)

TODO(among the easy to understand things that are tested, pick one short test and inline it here as an example)

The tests are located in the Kotlin repository at \texttt{compiler/testData/codegen/box/coroutines}

\section{Size benchmark}
\label{sec:size-bench}

\section{Performance benchmarks}
\label{sec:performance-bench}

\section{Summary}
\label{sec:evaluation-summary}

\chapter{Conclusion and Future Work}
\label{ch:conclusion}

\section{Summary}
\label{sec:summary}

\section{Main Contributions}
\label{sec:main-contributions}

\section{Key Findings}
\label{sec:key-findings}

\section{Limitations}
\label{sec:conclusion-limitations}

\section{Future Work}
\label{sec:future-work}

\section{Closing Remarks}
\label{sec:closing}


% Back matter
\backmatter
\bibliographystyle{plain}
\bibliography{references}

\appendix
\chapter{Stack-switching in WebAssembly}
\label{ch:appendix}

\section{Understanding Resuming and Suspending}
\label{sec:stack-switching-intuition}

A good intuition for understanding resuming and suspending in the stack-switching proposal is that it is similar to calling a function inside a try-catch block. The function can suspend by throwing different types of exceptions (suspension tags) and the try-catch block can provide different handlers for the corresponding suspension tags.

\section{WebAssembly Code Example}
\label{sec:wasm-example}

Below is a simplified example in Wasm that showcases the basics of using continuations (full example here). \texttt{\$consumer} creates a continuation object from the function \texttt{\$generator} and resumes it. When \texttt{\$generator} suspends with the tag \texttt{\$generate}, \texttt{\$consumer} jumps to the label \texttt{\$on\_generate}. At that point, the stack will contain parameters of the \texttt{\$generate} tag and a new continuation object. \texttt{\$consumer} saves the new continuation object and repeats the loop. Once \texttt{\$generator} finishes execution without suspending, \texttt{\$consumer} also returns.

\begin{lstlisting}[caption={WebAssembly stack-switching example}, label={lst:wasm-example}]
(func $generator
  ...
  (loop $loop
    ;; Suspend execution, pass current value of $i to consumer.
    (suspend $generate (local.get $i))
    ...
  )
)

(func $consumer
 ...
 ;; Create continuation executing function $generator.
 (local.set $generator_cont (cont.new $cont_type (ref.func $generator)))

 (loop $loop
   (block $on_generate (result i32 (ref $cont_type))
     ;; Resume continuation $generator_cont.
     ;; When $generate tag is thrown,
     ;; go to label $on_generate (points to right after this block)
     (resume $cont_type (on $generate $on_generate) (local.get $generator_cont))
     ;; $generator returned: no more data.
     (return)
   )
   ;; $generator suspended.
   ;; The stack now contains the value
   ;; produced by the generator and the new continuation.
   ;; Save continuation to resume it in the next iteration
   (local.set $generator_cont)
   ...
 )
)
\end{lstlisting}

\section{Control-Flow Graph}
\label{sec:control-flow-graph}

Figure~\ref{fig:stack-switching-cfg} illustrates the control flow between the \texttt{\$consumer} and \texttt{\$generator} functions, showing how stack-switching enables cooperative execution through continuations.

\begin{figure}[htbp]
\centering
\documentclass[tikz,border=10pt]{standalone}
\usepackage{tikz}
\usetikzlibrary{shapes,arrows,positioning,fit,calc}

\begin{document}
\begin{tikzpicture}[
    node distance=1.2cm and 3cm,
    every node/.style={font=\small},
    block/.style={rectangle, draw, fill=blue!10, text width=5em, text centered, rounded corners, minimum height=2.5em},
    decision/.style={diamond, draw, fill=yellow!10, text width=4em, text centered, aspect=2, inner sep=1pt},
    action/.style={rectangle, draw, fill=green!10, text width=6em, text centered, rounded corners, minimum height=2.5em},
    suspend/.style={rectangle, draw, fill=orange!20, text width=6em, text centered, rounded corners, minimum height=2.5em},
    arrow/.style={->, >=stealth, thick},
    transfer/.style={->, >=stealth, thick, dashed, color=red!70},
    solidarrow/.style={->, >=stealth, thick},
    label/.style={font=\scriptsize, text=red!70}
]

% Consumer side (left)
\node[block] (consumer-start) {\textbf{\$consumer} start};
\node[action, below=of consumer-start] (create-cont) {Create continuation from \$generator};
\node[block, below=of create-cont] (consumer-loop) {Loop start};
\node[suspend, below=of consumer-loop] (resume) {Resume continuation};
\node[block, below=of resume] (consumer-return) {Return};
\node[action, below=of consumer-return] (handle-result) {Handle result};
\node[block, below=of handle-result] (loop-end) {Loop end};

% Generator side (right)
\node[block, right=of consumer-start, xshift=2cm] (generator-start) {\textbf{\$generator} start};
\node[block, below=of generator-start] (generator-loop) {Loop start};
\node[action, below=of generator-loop] (compute) {Compute value};
\node[suspend, below=of compute] (suspend) {Suspend with value};
\node[block, below=of suspend] (loop-back-gen) {Loop end};
\node[block, below=of loop-back-gen] (generator-return) {Return};

% Consumer flow
\draw[arrow] (consumer-start) -- (create-cont);
\draw[arrow] (create-cont) -- (consumer-loop);
\draw[arrow] (consumer-loop) -- (resume);
%\draw[arrow] (resume) -- (handle-result);
\draw[arrow] (handle-result) -- (loop-end);
%\draw[arrow] (loop-end) -- (consumer-return);
\draw[arrow] (loop-end) to[out=180, in=180] (consumer-loop);

% Generator flow
\draw[arrow] (generator-start) -- (generator-loop);
\draw[arrow] (generator-loop) -- (compute);
\draw[arrow] (compute) -- (suspend);
%\draw[arrow] (suspend) -- (loop-back-gen);
\draw[arrow] (loop-back-gen) to[out=0, in=0] (generator-loop);
\draw[arrow] (loop-back-gen) -- node[above, pos=0.5] {no more values to produce} (generator-return);

% Stack-switching transfers
\draw[transfer] (create-cont) -- (generator-start);
\draw[arrow] (resume) -- node[above, pos=0.3] {first iteration} (generator-start);
\draw[arrow] (resume) -- (loop-back-gen);
\draw[arrow] (suspend) -- (handle-result);
\draw[arrow] (generator-return) -- (consumer-return);
%\draw[arrow] (handle-result) to[out=180, in=180] node[above, pos=0.5] {no value} (consumer-return);

\end{tikzpicture}

\end{document}

\caption{Control-flow graph showing stack-switching between \texttt{\$consumer} and \texttt{\$generator}. Solid arrows show sequential execution within each function. Dashed red arrows show control transfers via stack-switching operations: \texttt{cont.new} creates the continuation, \texttt{resume} transfers control to the generator, \texttt{suspend} transfers control back to the consumer with a value and updated continuation, and \texttt{return} signals completion.}
\label{fig:stack-switching-cfg}
\end{figure}


\end{document}

%     cont.new
%        |
%      resume -----> suspend
%        |  \
%        | handle
%      end